\subsection{Autorizzare una richiesta di acquisto}\label{sec:uc-autorizzazione}
\begin{xltabular}{\textwidth}{| >{\bfseries}p{.22\textwidth} | >{\raggedright\arraybackslash}X |}
\caption{Caso d'uso: Autorizzare una richiesta di acquisto}\label{tab:uc-autorizzazione}\\
\hline
\rowcolor{gray!15} Campo & \textbf{Contenuto}\\\hline
\endfirsthead
\multicolumn{2}{c}{{\tablename\ \thetable{} -- Continua dalla pagina precedente}}\\
\hline
\rowcolor{gray!15} Campo & \textbf{Contenuto}\\\hline
\endhead
\hline
\multicolumn{2}{r}{\textit{Continua nella pagina successiva...}}\\
\endfoot
\endlastfoot
Nome & Autorizzare una richiesta di acquisto\\\hline
Livello & Obiettivo Utente\\\hline
Descrizione & Il Salesman autorizza il Customer a completare un acquisto per cui è già stato versato un acconto. L'autorizzazione assegna il venditore alla pratica e stabilisce il termine per il saldo.\\\hline
Pre-condizioni & Il Salesman ha effettuato l'accesso alla propria area; esiste una richiesta di acquisto con acconto positivo e il relativo veicolo è stato riservato dal Customer.\\\hline
Post-condizioni & La richiesta conserva l'identificativo del Salesman e una scadenza pari alla data di autorizzazione più un mese. Il veicolo resta \code{BLOCKED}; in questa fase non avvengono addebiti.\\\hline
Attori & Salesman (primario), Customer (beneficiario dell'autorizzazione)\\\hline
Test & \code{CheckoutFunctionalTest.completePurchaseWorkflow} verifica che l'autorizzazione imposti la scadenza a un mese e che il successivo acquisto registri il venditore. La selezione nella GUI e i rami di errore non sono verificati da quel test.\\\hline
Flusso Principale &
\begin{enumerate}[label=\arabic*.,start=1,leftmargin=*,nosep]
  \item Il Salesman consulta le richieste di acquisto e seleziona la pratica da autorizzare.
  \item Attiva il comando di autorizzazione.
  \item Il sistema recupera la richiesta e verifica che si tratti di un acquisto, riconoscibile dall'acconto positivo.
  \item Il sistema salva nella richiesta l'identificativo del Salesman e la data di scadenza, in un'unica transazione.
  \item L'interfaccia conferma l'esito e aggiorna l'elenco. Il Customer può procedere al saldo dalla propria area.
\end{enumerate}\\\hline
Flussi Alternativi &
\textbf{A1 (passo 1). Nessuna richiesta selezionata:} l'interfaccia mostra un messaggio e non invia l'autorizzazione.\newline\newline
\textbf{A2 (passo 3). Richiesta assente o di tipo diverso:} una richiesta non più presente produce un errore; una proposta di cessione è rifiutata dall'interfaccia e dal servizio di autorizzazione. Nessuna autorizzazione viene salvata.\newline\newline
\textbf{A3 (passo 4). Salvataggio non riuscito:} la transazione viene annullata e l'interfaccia comunica l'errore.\\\hline
\end{xltabular}
